\chapter{Indigestion}
\index{Indigestion}

\section*{Definition and Overview}
Indigestion, also known as dyspepsia, is a term for a group of symptoms that arise from the upper part of the digestive tract, including discomfort or pain in the upper abdomen, feeling full quickly during a meal, bloating, belching, and nausea. It is extremely common and usually short-lived, often related to eating habits, stress, or medicines. In many people the cause is not found on investigation, a situation known as functional dyspepsia. In others, indigestion results from conditions such as gastro-oesophageal reflux, inflammation of the stomach lining (gastritis), or a stomach ulcer. Rarely, it signals something more serious, so persistent or severe symptoms should always be assessed by a doctor.

\section*{Epidemiology}
Indigestion affects around one in four people at some time, and up to one in five people experience it regularly. It is more common with increasing age, in people who smoke, and in those taking medicines that irritate the stomach, such as non-steroidal anti-inflammatory drugs (NSAIDs). Rates are higher in people who are overweight, drink alcohol, or have chronic stress. Most people with indigestion never seek medical care, and of those who do, most are found to have no serious underlying disease. The condition has a significant effect on quality of life, work, and sleep.

\section*{Aetiology and Pathophysiology}
Indigestion arises from irritation or dysfunction of the upper gastrointestinal tract. Stomach acid is a normal part of digestion, but when the protective lining of the stomach is impaired, or when acid refluxes into the oesophagus, symptoms develop. Common causes include overeating, eating too quickly, fatty or spicy foods, caffeine, alcohol, and smoking. Medicines such as NSAIDs and some antibiotics can damage the stomach lining. Helicobacter pylori infection causes chronic inflammation of the stomach lining and is a major cause of ulcers. In functional dyspepsia, the stomach empties slowly or is unusually sensitive to distension, without any visible abnormality. Stress and anxiety amplify the perception of symptoms through the gut-brain connection.

\section*{Clinical Features}
\begin{itemize}
    \item Pain, burning, or discomfort in the upper abdomen
    \item Feeling full very quickly during meals or uncomfortably full afterwards
    \item Bloating, belching, and excessive wind
    \item Nausea, sometimes with retching
    \item Heartburn: a burning sensation rising from the stomach into the chest
    \item Symptoms often related to eating, stress, or certain foods
    \item Symptoms may occur at night and disturb sleep
\end{itemize}

\section*{Differential Diagnosis}
Indigestion can be confused with other conditions that cause upper abdominal or chest symptoms.
\begin{itemize}
    \item \textbf{Gastro-oesophageal reflux disease:} acid rising into the oesophagus
    \item \textbf{Gastritis and peptic ulcers:} inflammation or ulceration of the stomach or duodenum
    \item \textbf{Gallstones:} pain after fatty meals, usually in the right upper abdomen
    \item \textbf{Gastroparesis:} slow stomach emptying, common in diabetes
    \item \textbf{Coeliac disease:} gluten intolerance causing bloating and nausea
    \item \textbf{Stomach cancer:} rare, but suspected with weight loss, anaemia, or difficulty swallowing
    \item \textbf{Pancreatic disease:} pain radiating to the back
    \item \textbf{Angina or heart attack:} chest discomfort that must be excluded, especially in older people and those with heart disease
\end{itemize}

\section*{When to Seek Medical Care}
See a doctor if indigestion is severe, persistent, or worsening; if it occurs with weight loss, difficulty swallowing, vomiting, or black stools; or if it began after age 55. People taking NSAIDs or aspirin regularly, and those with a family history of stomach cancer, should also be assessed. Review medicines with a doctor or pharmacist, as many cause indigestion.

\textbf{Contact your local emergency services for:}
\begin{itemize}
    \item Chest pain, pressure, or tightness, especially with breathlessness, sweating, or pain in the arm or jaw
    \item Vomiting blood or material that looks like coffee grounds
    \item Black, tarry, or bloody stools
    \item Sudden severe abdominal pain
    \item Difficulty swallowing or choking
\end{itemize}

\section*{Diagnosis and Investigations}
\begin{itemize}
    \item \textbf{History:} timing and triggers of symptoms, diet, medicines, smoking, and alcohol
    \item \textbf{Examination:} abdominal examination and assessment for anaemia or weight loss
    \item \textbf{Helicobacter pylori testing:} a breath test, stool test, or biopsy to detect the infection
    \item \textbf{Gastroscopy (endoscopy):} direct examination of the oesophagus, stomach, and duodenum with biopsies, especially in older people or those with alarm symptoms
    \item \textbf{Trial of acid-suppressing treatment:} a short course of a proton pump inhibitor to see if symptoms resolve
    \item \textbf{Blood tests:} full blood count for anaemia, and thyroid, liver, or pancreatic tests as indicated
    \item \textbf{Imaging:} ultrasound for gallstones or other upper abdominal problems in selected cases
\end{itemize}

\section*{Management}
\begin{itemize}
    \item \textbf{Lifestyle measures:} eating smaller, more frequent meals; avoiding trigger foods; limiting alcohol and caffeine; and stopping smoking
    \item \textbf{Weight management:} losing weight if overweight, as this reduces pressure on the stomach
    \item \textbf{Medicines:} antacids for quick relief, and acid-suppressing drugs such as proton pump inhibitors or H2 blockers
    \item \textbf{Treating H. pylori:} a course of antibiotics and acid suppression when the infection is found
    \item \textbf{Reviewing medicines:} stopping or changing NSAIDs and other culprits where possible
    \item \textbf{Stress management:} relaxation, mindfulness, and addressing anxiety
    \item \textbf{Treating underlying conditions:} reflux, ulcers, gallstones, or gastroparesis
    \item \textbf{Functional dyspepsia:} may respond to low-dose antidepressants or prokinetic medicines that improve stomach emptying
\end{itemize}

\section*{Complications}
\begin{itemize}
    \item Peptic ulcers with bleeding, perforation, or obstruction
    \item Anaemia from chronic blood loss
    \item Oesophageal damage and strictures from long-standing reflux
    \item Barrett's oesophagus, a precursor to oesophageal cancer in severe reflux
    \item Weight loss and malnutrition from fear of eating
    \item Reduced quality of life, sleep disturbance, and anxiety
    \item Delayed diagnosis of more serious conditions when symptoms are dismissed
\end{itemize}

\section*{Prognosis and Outcomes}
Most episodes of indigestion settle with simple lifestyle changes and over-the-counter remedies. When H. pylori infection is treated, ulcers heal and most stay healed. Functional dyspepsia tends to run a relapsing course, but symptoms are manageable in most people with a combination of lifestyle measures, stress management, and medication. Investigation provides reassurance and identifies the small number of people with serious disease, in whom early treatment dramatically improves outcomes.

\section*{Prevention and Self-Care}
\begin{itemize}
    \item Eat smaller, slower meals and avoid late-night eating
    \item Identify and avoid personal trigger foods: fatty, spicy, or acidic foods
    \item Limit alcohol and caffeine, and stop smoking
    \item Maintain a healthy weight
    \item Avoid NSAIDs unless advised, and take them with food if needed
    \item Avoid tight clothing and lying down soon after meals
    \item Manage stress with exercise, relaxation, and good sleep
    \item See a doctor for persistent or severe symptoms rather than self-treating long term
    \item Review regular medicines with a doctor or pharmacist
\end{itemize}

\section*{Sources and Further Reading}
The following reputable sources provide further information for healthcare students and patients seeking reliable, current advice about indigestion.
\begin{itemize}
    \item Healthdirect Australia. \textit{Indigestion (dyspepsia).} https://www.healthdirect.gov.au/
    \item Gastroenterological Society of Australia. \textit{Patient information on dyspepsia.} https://www.gesa.org.au/
    \item NHS. \textit{Indigestion.} https://www.nhs.uk/conditions/indigestion/
    \item Mayo Clinic. \textit{Indigestion.} https://www.mayoclinic.org/
    \item MSD Manual. \textit{Dyspepsia.} https://www.msdmanuals.com/
    \item Better Health Channel. \textit{Indigestion and heartburn.} https://www.betterhealth.vic.gov.au/
\end{itemize}
